\begin{center}
\textsc{Gauss} an \textsc{Gerling}.\quad Göttingen, 3.~October 1823.
\end{center}

\par\noindent\hspace*{\parindent}\mbox{.\ .\ .\ .\ .\ .}\quad Erlauben Sie mir noch einige \mbox{Bemerkungen.\quad\mbox{.\ .\ .\ .\ .\ .}}
\par\noindent\hspace*{\parindent}1) Beim Vergleichen zweier Punkte fange ich jedesmal von demjenigen
an, von dem ich am meisten befürchte, dass er versagen könnte.
\par\clearpage
\noindent\makebox[0pt][l]{\fontsize{13}{16}\selectfont 390}\hfill{\small\textbf{BRIEFWECHSEL MIT GERLING.}}\hfill\null\par\medskip
\par\noindent\hspace*{\parindent}2) Ist der zweite Punkt Heliotroplicht und nicht die allergrösste Hoffnung,
dass er nicht versagen wird, so lese ich öfters ab als gewöhnlich. Wo diese
Furcht nur äusserst gering war, habe ich wohl zuweilen mehr als 10, d.~i.\ 15,
ja wohl 20 Messungen gemacht, ehe ich wieder ablas. Inzwischen setzt man
dabei immer viel aufs Spiel. Ist viel Besorgniss des Versagens des zweiten
Punktes, so lese ich wohl nach der dritten, ja nach der zweiten oder gar bei
jeder Beobachtung ab. Dasselbe auch, wenn der erste Punkt etwas lange
ausbleibt.
\par\noindent\hspace*{\parindent}3) Bin ich in einer Beobachtungsreihe begriffen, wo der erste Punkt
geschnitten ist, so lasse ich gern, ehe ich die Alhidade löse, durch ein
Handfernrohr nachsehen, ob der zweite Punkt sichtbar ist; versagt er aber,
nachdem schon gelöst ist und kommt nicht wieder, so ist deswegen doch die
Messung nicht verloren; ich schneide dann im Nothfall einen dritten Punkt ein
und bekomme so gemischte Winkel. Bei meinem diesmaligen Brockenbesuch
ist indessen dieser Fall nur Einmal eingetreten, indem ich 2-mal
(Ilefeld-Hohehagen) \(+\) 2-mal (Ilefeld-Hils) nahm. Solche gemischte Winkel werden im
System ebenso gut und ganz pro rata benutzt wie reine. Im Jahr 1821 habe
ich viele gemischte Winkel, später seltener. Ich habe daher auch an den
meisten Stationen, ausser den Hauptrichtungen, Hülfsrichtungen; besonders im
Jahre 1821, wo ich nur Einen Heliotrop hatte. Ich wähle zu Hülfsrichtungen
gern spitze Thürme, nicht gar zu entfernt und wo möglich nicht gar zu arg
ausser der Horizontalebene. Letztere Bedingung war freilich auf dem Brocken
nicht zu erreichen. Huyseburg liegt \(1^{0}\,37^{\prime}\), Hüttenrode \(1^{0}\,54^{\prime}\) und eine Stange
auf dem Wurmberg, die ich 1821 aufpflanzen und 1823 erneuern liess
(obwohl an einem etwas andern Platz), \(2^{0}\,13^{\prime}\) unter dem Horizont.
\par
